\chapter{Configuración de entrenamiento de los modelos}
\label{anexo:configuracion-entrenamiento}

Este anexo reúne la configuración con la que se entrenaron los dos modelos de segmentación, con el detalle necesario para que un tercero pueda reproducir los resultados informados en el Capítulo 13. Se incluyen los parámetros de optimización, la estrategia de aumentación de datos y los controles de calidad aplicados al cierre del entrenamiento. % CROSS_REFERENCE_MAPPING_REQUIRED: Capítulo 13

\phantomsection
\section*{E.1 Parámetros de entrenamiento}
\addcontentsline{toc}{section}{E.1 Parámetros de entrenamiento}
\label{anexo:entrenamiento-f1-parametros}

\begin{table}[H]
\small
\centering
\caption{Configuración comparada del entrenamiento de ambos modelos.}
\label{tab:anexo-e-parametros-entrenamiento}
\begin{tabularx}{\textwidth}{|>{\raggedright\arraybackslash}p{0.25\textwidth}|>{\raggedright\arraybackslash}X|>{\raggedright\arraybackslash}X|}
\hline
\textbf{Parámetro} & \textbf{Modelo sagital} & \textbf{Modelo axial} \\
\hline
Conjunto de datos & SPIDER & Sudirman y Al-Kafri \\
\hline
Arquitectura & U-Net bidimensional, variante sagital & U-Net bidimensional, variante axial \\
\hline
Número de clases & 4 & 6 \\
\hline
Canales base & 16 & 16 \\
\hline
Dimensión de entrada & $256 \times 256$ & $256 \times 256$ \\
\hline
Semilla & 20260716 & 2026 \\
\hline
Tamaño de lote & 8 & 8 \\
\hline
Épocas máximas & 80 & 80 \\
\hline
Tasa de aprendizaje & $8 \times 10^{-4}$ & $8 \times 10^{-4}$ \\
\hline
Decaimiento de pesos & $1 \times 10^{-4}$ & $1 \times 10^{-4}$ \\
\hline
Paciencia de detención temprana & 12 & 12 \\
\hline
Paciencia del planificador & 4 & No aplicada \\
\hline
Recorte de gradiente & No aplicado & 1,0 \\
\hline
Precisión mixta & Sí & Sí \\
\hline
Función de pérdida & 0,5 $\times$ entropía cruzada ponderada + 0,5 $\times$ Dice sin fondo & Entropía cruzada ponderada + Dice suave sin fondo \\
\hline
Pesos de clase & Estimados del conjunto & Estimados del conjunto, con refuerzo sobre la clase de fondo de origen \\
\hline
Épocas completadas & 75, con detención temprana & Detención temprana \\
\hline
Época seleccionada & 63 & 17 \\
\hline
\end{tabularx}
\end{table}

La función de pérdida combina en ambos casos un término de entropía cruzada ponderada por clase con un término basado en el coeficiente Dice que excluye el fondo. La combinación responde a que la entropía cruzada por sí sola tiende a favorecer las clases mayoritarias, mientras que el término de Dice equilibra la contribución de las estructuras de menor tamaño.

\phantomsection
\section*{E.2 Aumentación de datos}
\addcontentsline{toc}{section}{E.2 Aumentación de datos}
\label{anexo:entrenamiento-f2-aumentacion}

La estrategia de aumentación difiere entre ambos modelos, y la diferencia responde a una razón anatómica que conviene explicitar.

\begin{table}[H]
\small
\centering
\caption{Transformaciones de aumentación aplicadas durante el entrenamiento.}
\label{tab:anexo-e-aumentacion}
\begin{tabularx}{\textwidth}{|>{\raggedright\arraybackslash}p{0.25\textwidth}|>{\raggedright\arraybackslash}X|>{\raggedright\arraybackslash}X|}
\hline
\textbf{Transformación} & \textbf{Modelo sagital} & \textbf{Modelo axial} \\
\hline
Reflexión horizontal & Habilitada, probabilidad 0,50 & Deshabilitada \\
\hline
Reflexión vertical & Habilitada, probabilidad 0,20 & Deshabilitada \\
\hline
Rotación & Habilitada, $\pm 8$ grados, probabilidad 0,40 & Deshabilitada \\
\hline
Variación de contraste & Habilitada, factor entre 0,85 y 1,15, probabilidad 0,50 & Habilitada, factor entre 0,95 y 1,05, probabilidad 0,25 \\
\hline
Interpolación & Bilineal para la imagen, vecino más próximo para la máscara & Idéntica \\
\hline
\end{tabularx}
\end{table}

La reflexión horizontal está habilitada en el modelo sagital y deshabilitada en el axial. En un corte axial, el lado izquierdo y el derecho corresponden a regiones anatómicas distintas y clínicamente distinguibles, de modo que reflejar la imagen invierte la lateralidad de cualquier hallazgo. En un corte sagital no existe lateralidad que preservar. Habilitar la reflexión en el plano axial habría multiplicado los datos disponibles a costa de entrenar el modelo para que ignore una distinción que el uso clínico requiere.

La rotación permanece deshabilitada en el módulo axial por un motivo análogo: la angulación de un corte axial es un dato del estudio, no una variación arbitraria, y la asignación de nivel anatómico descrita en la sección 11.6 depende de ella. % CROSS_REFERENCE_MAPPING_REQUIRED: sección 11.6

Ambas estrategias son deterministas por época. La semilla de cada transformación se deriva de la combinación de la semilla global, el número de época, el índice de la muestra y el identificador del corte, de modo que la reproducción es exacta corte por corte y no únicamente a nivel de ejecución completa.

\phantomsection
\section*{E.3 Controles de calidad al cierre del entrenamiento}
\addcontentsline{toc}{section}{E.3 Controles de calidad al cierre del entrenamiento}
\label{anexo:entrenamiento-f3-controles-calidad}

\begin{table}[H]
\small
\centering
\caption{Condiciones evaluadas antes de promover un artifact.}
\label{tab:anexo-e-controles-promocion-artifact}
\begin{tabularx}{\textwidth}{|>{\raggedright\arraybackslash}p{0.40\textwidth}|>{\raggedright\arraybackslash}X|>{\raggedright\arraybackslash}X|}
\hline
\textbf{Control} & \textbf{Valor o condición} & \textbf{Configurable} \\
\hline
Umbral de coeficiente Dice macro de primer plano & 0,70 & Sí \\
\hline
Partición retenida separada a nivel de paciente & Obligatoria & No \\
\hline
Informe explícito del desempeño sobre la clase de fondo de origen & Obligatorio & No \\
\hline
Informe de intersección sobre unión además del coeficiente Dice & Obligatorio & No \\
\hline
Generación de una partición nueva & Deshabilitada por omisión & Requiere habilitación manual \\
\hline
\end{tabularx}
\end{table}

El cuaderno de entrenamiento incorpora además verificaciones de fuga cruzada entre particiones, de particiones vacías y de clases ausentes en alguna de ellas, así como una restricción en tiempo de ejecución que impide reevaluar la partición retenida de una ejecución ya evaluada sin confirmación explícita.

Según se informa en la sección 13.4, el umbral no se alcanzó en el modelo axial bajo su lectura de primer plano completo, y el artifact quedó registrado en esa condición sin que el valor del umbral se modificara. % CROSS_REFERENCE_MAPPING_REQUIRED: sección 13.4

\phantomsection
\section*{E.4 Auditoría documental de la partición axial E9}
\addcontentsline{toc}{section}{E.4 Auditoría documental de la partición axial E9}
\label{anexo:entrenamiento-e4-auditoria-e9}

La partición E9 del subconjunto axial curado reúne 610 registros: 427 de entrenamiento (128 identificadores de paciente), 81 de validación (27) y 102 de prueba (29), para 184 identificadores de paciente sin intersecciones entre particiones.

El CSV auditado es \path{E9_t2_final_labels_curated_split.csv}. Su SHA-256 coincide con el registrado en el manifest del artifact axial V2:

{\small\ttfamily 17f78265\allowbreak b93cca59\allowbreak b0680c5f\allowbreak e336943a\allowbreak cee78c4f\allowbreak be34bd86\allowbreak a23df06c\allowbreak e80d051e}.

La auditoría documental H11-A reportó cero filas idénticas duplicadas y ninguna intersección de identificadores de pacientes entre entrenamiento, validación y prueba. Una verificación posterior H11-B, basada en los archivos referenciados por el CSV, resolvió 610 de 610 rutas exactas de imagen y máscara y calculó hashes de contenido para 1.162 archivos distintos, sin errores de lectura. No encontró grupos de imágenes idénticas ni de pares imagen-máscara idénticos entre particiones. Encontró 55 grupos de máscaras idénticas asociados con 113 filas; todos permanecen dentro de un mismo split y paciente, sin grupos entre splits ni pacientes. La repetición de una máscara —incluida una máscara vacía— no demuestra por sí sola fuga ni correspondencia anatómica correcta.

La verificación H11-B no reemplaza la correspondencia semántica de cada máscara con su imagen, la geometría/anatomía, ni la auditoría de metadatos de paciente y serie. Tampoco se comparó con la otra partición histórica porque no se dispuso del segundo CSV. Los identificadores de estudio informados derivan de \texttt{case\_id\_norm}, igual que los identificadores de paciente, por lo que no prueban de manera independiente el número de estudios clínicos distintos. La igualdad SHA-256 prueba identidad de bytes, no alineación anatómica ni corrección de etiquetas.

La partición de prueba ya había sido utilizada con \texttt{axial-full-v1}; el resultado V2 es comparativo y no constituye una evaluación sobre un test completamente intacto. En consecuencia, el hallazgo \#11 se registra como parcialmente resuelto: los conteos, el hash y la separación por identificadores están documentados, mientras que los controles anteriores permanecen pendientes o son una limitación confirmada. Esta incorporación no modifica las métricas ni el estado de candidato del artifact.

El handoff de H11-A, junto con el resumen y reporte JSON anonimizado de H11-B, se conserva en el expediente local de evidencias bajo \path{docs/audit/pfi_100_tutor_2026_10_14/evidencias/fuentes_francisco_2026-10-14/}. El artifact axial sigue siendo candidato; no se modificaron métricas, manifests ni gates.
\phantomsection
\section*{E.5 Trazabilidad de los artifacts}
\addcontentsline{toc}{section}{E.5 Trazabilidad de los artifacts}
\label{anexo:entrenamiento-f4-trazabilidad-artifacts}

Cada artifact final conserva en su manifest la correspondencia de clases, la configuración de la arquitectura, la dimensión de entrada, las métricas obtenidas, el estado de entrenamiento y un conjunto de identificadores criptográficos que vinculan el modelo con el código y los datos que lo produjeron: el resumen del índice de entrenamiento, el del código de la arquitectura, el commit del repositorio y la referencia al cuaderno de origen.

El runtime valida esa correspondencia antes de habilitar un modelo. Si los nombres de clase declarados en el código no coinciden con los del manifest, el modelo queda deshabilitado en lugar de ejecutarse con una interpretación distinta de sus salidas.

\phantomsection
\section*{E.6 Desglose histórico sagital recuperado}
\addcontentsline{toc}{section}{E.6 Desglose histórico sagital recuperado}
\label{anexo:entrenamiento-e6-h04-historico}

El paquete anonimizado de auditoría H04 contiene un inventario de un archivo histórico llamado \texttt{final\_test\_metrics.json}, con las métricas resumidas en la Tabla~\ref{tab:anexo-e-h04-metricas-historicas}. El SHA-256 del archivo original que informa ese inventario es \texttt{0b23514d\allowbreak 88f75436\allowbreak 154fc66c\allowbreak bcf1e62a\allowbreak 71c2b220\allowbreak 86624473\allowbreak 44e081b2\allowbreak d344818eb}; el archivo completo no se encuentra en el paquete ni en el snapshot local inspeccionado, por lo que no fue posible recalcular ese hash. En consecuencia, las cifras son evidencia histórica secundaria pendiente de autenticación binaria y no representan una nueva inferencia ni un recálculo efectuado durante esta auditoría.

\begin{table}[H]
\small
\centering
\caption{Métricas por clase transcritas del inventario anonimizado del reporte histórico sagital.}
\label{tab:anexo-e-h04-metricas-historicas}
\begin{tabularx}{\textwidth}{|>{\raggedright\arraybackslash}p{0.23\textwidth}|>{\centering\arraybackslash}p{0.18\textwidth}|>{\centering\arraybackslash}p{0.18\textwidth}|>{\centering\arraybackslash}p{0.19\textwidth}|>{\centering\arraybackslash}X|}
\hline
\textbf{Clase} & \textbf{Dice} & \textbf{IoU} & \textbf{Cortes GT} & \textbf{Cortes con predicción} \\
\hline
Fondo (excluido del macro) & 0,993858 & 0,987790 & 1.218 & 1.218 \\
\hline
Grupo vertebral & 0,908631 & 0,832561 & 1.113 & 1.107 \\
\hline
Canal & 0,907121 & 0,830029 & 465 & 483 \\
\hline
Grupo discal & 0,864543 & 0,761405 & 887 & 885 \\
\hline
Macro de primer plano (clases 1--3) & 0,893432 & 0,807998 & -- & -- \\
\hline
\end{tabularx}
\end{table}

Los valores Dice e IoU de la tabla se redondean a seis decimales para facilitar la lectura; la precisión completa permanece en el JSON anonimizado archivado.

El inventario atribuye el reporte a 1.218 cortes de $256 \times 256$ y a 33 pacientes de prueba según el manifest; este cuaderno no vuelve a contar pacientes. El macro es aritmético entre las tres clases de primer plano, después de agregar por clase sobre píxeles/cortes; no es un macro por paciente ni un resultado del benchmark oculto SPIDER. El hash del índice de entrenamiento consignado es \texttt{2720b721\allowbreak 8c92870f\allowbreak 6f0a000b\allowbreak 57111ed3\allowbreak 6b5cf3b7\allowbreak 8c716f24\allowbreak 4f427ca7\allowbreak fee4a4ba}, coincidente con el manifest y el cuaderno de entrenamiento inspeccionados. Sin embargo, el snapshot contiene el checkpoint runtime de SHA-256 \texttt{cf11dcc0\allowbreak ad77a7c7\allowbreak 87e64a79\allowbreak 6a2fd739\allowbreak 8ef906ad\allowbreak d461cef4\allowbreak b3d61f1a\allowbreak 5238e944}, mientras que el inventario informa un best checkpoint de entrenamiento con SHA-256 \texttt{fb124f0b\allowbreak 1e73a5f3\allowbreak e548769b\allowbreak f8370892\allowbreak 95ce8e02\allowbreak c33fb34e\allowbreak d3be7fc3\allowbreak c2142a39}. No se estableció equivalencia de tensores, mapping, arquitectura y protocolo entre ambos; no se infiere esa equivalencia a partir de la coincidencia de métricas.

El paquete anonimizado conserva además el resultado \texttt{NO\_RECALCULADO}: no contenía las predicciones/arrays ni el checkpoint original requeridos para repetir la evaluación. El ZIP, el cuaderno H04 y el handoff recibido se archivan en \path{docs/audit/pfi_100_tutor_2026-10-14/evidencias/fuentes_h04_2026-10-14/}. H04 permanece parcial hasta localizar el JSON original y resolver la procedencia del checkpoint y del protocolo.
